\chapter{Grafos II: Caminhos e Árvores}

\section{Dijkstra}

\subsection{Padrão do Problema}
Caminhos mínimos a partir de uma única origem em grafo com pesos não negativos.

\subsection{Quando usar}
\begin{itemize}[nosep]
  \item Todos os pesos das arestas são $\geq 0$.
  \item É preciso a menor distância de uma origem para vários destinos.
  \item O grafo é grande demais para Floyd-Warshall.
\end{itemize}

\subsection{Explicação do Algoritmo}
Mantenha distâncias tentativas e uma fila de prioridade com o vértice de menor distância atual. Ao processar um vértice, relaxe suas arestas.

\subsection{Complexidade Temporal}
$O((V + E) \log V)$ com heap binário.

\subsection{Exemplo de Enunciado}
\begin{exemplodiabox}
\textbf{Enunciado:} Dado um mapa de cidades com distâncias não negativas, calcule a menor distância da cidade $s$ para todas as outras.
\end{exemplodiabox}

\subsection{Código}
\incluircppcaption{algoritmos_avulsos/grafos/dijkstra.cpp}{Dijkstra}

\section{Bellman-Ford}

\subsection{Padrão do Problema}
Caminhos mínimos com possibilidade de arestas negativas e detecção de ciclos negativos.

\subsection{Quando usar}
\begin{itemize}[nosep]
  \item Existem pesos negativos.
  \item É necessário detectar ciclo negativo alcançável.
  \item O tamanho do grafo permite $O(VE)$.
\end{itemize}

\subsection{Explicação do Algoritmo}
Relaxe todas as arestas por $V-1$ rodadas. Uma rodada extra que ainda melhora alguma distância indica ciclo negativo.

\subsection{Complexidade Temporal}
$O(VE)$.

\subsection{Exemplo de Enunciado}
\begin{exemplodiabox}
\textbf{Enunciado:} Em uma rede de trocas com custos positivos e negativos, determine menores custos a partir da origem e detecte se há ciclo de ganho infinito.
\end{exemplodiabox}

\subsection{Código}
\incluircppcaption{algoritmos_avulsos/grafos/bellman_ford.cpp}{Bellman-Ford}

\section{Floyd-Warshall}

\subsection{Padrão do Problema}
Caminhos mínimos entre todos os pares de vértices em grafo pequeno.

\subsection{Quando usar}
\begin{itemize}[nosep]
  \item É preciso responder distância entre qualquer par.
  \item $V$ é pequeno o suficiente para $O(V^3)$.
  \item Pode haver peso negativo, mas não ciclo negativo.
\end{itemize}

\subsection{Explicação do Algoritmo}
Use \texttt{dist[i][j]} e tente melhorar cada par usando um vértice intermediário $k$. A cada rodada, caminhos podem passar pelos vértices já liberados.

\subsection{Complexidade Temporal}
$O(V^3)$.

\subsection{Exemplo de Enunciado}
\begin{exemplodiabox}
\textbf{Enunciado:} Dadas até 400 cidades, responda a menor distância entre qualquer par de cidades.
\end{exemplodiabox}

\subsection{Código}
\incluircppcaption{algoritmos_avulsos/grafos/floyd_warshall.cpp}{Floyd-Warshall}

\section{Kruskal}

\subsection{Padrão do Problema}
Conectar todos os vértices com custo total mínimo, escolhendo arestas sem formar ciclo.

\subsection{Quando usar}
\begin{itemize}[nosep]
  \item O problema pede árvore geradora mínima.
  \item As arestas estão disponíveis como lista.
  \item O grafo é esparso ou ordenar arestas é conveniente.
\end{itemize}

\subsection{Explicação do Algoritmo}
Ordene as arestas por peso. Percorra da menor para a maior e adicione a aresta se ela conecta componentes diferentes, usando DSU.

\subsection{Complexidade Temporal}
$O(E \log E)$.

\subsection{Exemplo de Enunciado}
\begin{exemplodiabox}
\textbf{Enunciado:} Conecte todas as cidades com cabos de menor custo total, sem precisar manter ciclos.
\end{exemplodiabox}

\subsection{Código}
\incluircppcaption{algoritmos_avulsos/grafos/kruskal.cpp}{Kruskal}

\section{Prim}

\subsection{Padrão do Problema}
Construir uma árvore geradora mínima expandindo uma componente a partir de um vértice inicial.

\subsection{Quando usar}
\begin{itemize}[nosep]
  \item O problema pede conectar todos os vértices com custo mínimo.
  \item O grafo está em lista de adjacência ponderada.
  \item É natural crescer a solução a partir de um vértice.
\end{itemize}

\subsection{Explicação do Algoritmo}
Mantenha uma fila de prioridade com arestas que saem da árvore atual. Sempre escolha a menor aresta que leva a um vértice ainda fora da árvore.

\subsection{Complexidade Temporal}
$O((V + E) \log V)$ com heap.

\subsection{Exemplo de Enunciado}
\begin{exemplodiabox}
\textbf{Enunciado:} A partir de uma estação inicial, conecte todos os pontos de uma rede usando custo total mínimo.
\end{exemplodiabox}

\subsection{Código}
\incluircppcaption{algoritmos_avulsos/grafos/prim.cpp}{Prim}

\section{Union-Find (DSU)}

\subsection{Padrão do Problema}
Manter componentes disjuntas sob operações de união e consulta de conectividade.

\subsection{Quando usar}
\begin{itemize}[nosep]
  \item É preciso saber se dois elementos estão no mesmo grupo.
  \item Arestas são adicionadas dinamicamente.
  \item Algoritmos como Kruskal precisam detectar ciclos.
\end{itemize}

\subsection{Explicação do Algoritmo}
Cada conjunto tem um representante. \texttt{find} encontra esse representante com compressão de caminho; \texttt{unite} junta dois conjuntos por tamanho ou rank.

\subsection{Complexidade Temporal}
$O(\alpha(n))$ amortizado por operação, praticamente constante.

\subsection{Exemplo de Enunciado}
\begin{exemplodiabox}
\textbf{Enunciado:} Receba uniões entre pessoas e responda se duas pessoas já pertencem ao mesmo grupo.
\end{exemplodiabox}

\subsection{Código}
\incluircppcaption{algoritmos_avulsos/grafos/union_find.cpp}{Union-Find}

\section{LCA}

\subsection{Padrão do Problema}
Consultas de ancestral comum mais baixo em uma árvore enraizada.

\subsection{Quando usar}
\begin{itemize}[nosep]
  \item É preciso responder muitas consultas entre pares de nós.
  \item O problema envolve distância ou caminho em árvore.
  \item A árvore é fixa após o pré-processamento.
\end{itemize}

\subsection{Explicação do Algoritmo}
Com binary lifting, pré-calcule ancestrais em potências de dois. Para consultar, suba o nó mais profundo e depois suba ambos até antes do LCA.

\subsection{Complexidade Temporal}
Pré-processamento $O(n \log n)$ e consulta $O(\log n)$.

\subsection{Exemplo de Enunciado}
\begin{exemplodiabox}
\textbf{Enunciado:} Em uma árvore genealógica, responda várias consultas sobre o ancestral comum mais baixo de dois nós.
\end{exemplodiabox}

\subsection{Código}
\incluircppcaption{algoritmos_avulsos/grafos/lca.cpp}{LCA}
